\section{Esperimento con gli Utenti}

\subsection{Tipologia di utente presa in esame}
Si è deciso di concentrare l'esperimento esclusivamente sulla componente \textit{client} dell'applicazione, in quanto la maggior parte degli utenti finali interagirà prevalentemente con questa view. Tale scelta è dettata anche dalla natura fittizia del contesto applicativo: non disponendo di utenti reali con ruoli specifici di \textit{admin} o \textit{employee}, non è stato possibile coinvolgere profili professionali di questo tipo nella fase di test.

\subsection{Distribuzione e numeri degli utenti}
Per quanto riguarda gli utenti clienti della banca è stato facile recuperarli in quanto il target di riferimento coincide con la popolazione generale, trattandosi di un servizio potenzialmente utilizzato da qualsiasi individuo che faccia uso di applicazioni bancarie.
\subsubsection{Distribuzione utenti per sesso}
\subsubsection{Distribuzione utenti per età}
\subsubsection{Distribuzione utenti per grado di istruzione}

\subsection{Obiettivo}
L'obiettivo di questo esperimento è quello di confrontare l'applicazione bancaria prima della riprogettazione con l'applicazione dopo riprogettazione 

\subsection{Condizioni sperimentali}
Inserire descrizione delle variabili di controllo, indipendenti e dipendenti + Ambiente

\subsection{Approccio}
Scegliere tra whitin/between

\subsection{Progettazione temporale delle attività}

\subsection{Compiti}
Scegliere tra lista di operazioni o descrizione 

\subsection{Documentazione utilizzata}
